\chapter[2002 --- Brenner et al.]{2002: Sydney Brenner, H. Robert Horvitz, John E. Sulston}
\label{ch:2002_Brenner}

\section*{Main Contribution}

\textit{Discovered the genetic regulation of organ development and programmed cell death (apoptosis), explaining how the body sculpts its structures and eliminates unneeded cells.}\cite{nobel-2002,lecture-2002}

\section{Official Citation}

for their discoveries concerning genetic regulation of organ development and programmed cell death

\section{Background and Historical Context}

Programmed cell death---apoptosis---is as essential to the development and maintenance of multicellular organisms as cell division. During embryonic development, cells are produced in excess and those that are unnecessary or potentially dangerous are eliminated by a genetically controlled self-destruct program. The fingers and toes of the developing hand, for example, are sculpted by the death of the webbing between them. The nervous system is sculpted by the death of neurons that fail to make appropriate synaptic connections. In the adult body, apoptosis eliminates damaged, infected, or potentially cancerous cells. The genetic mechanisms controlling programmed cell death were elucidated through the work of Sydney Brenner, H. Robert Horvitz, and John E. Sulston, whose discoveries earned them the 2002 Nobel Prize in Physiology or Medicine.

Sydney Brenner was born on January 13, 1928, in Germiston, South Africa. He studied medicine at the University of the Witwatersrand in Johannesburg and then earned his D.Phil.\ from Oxford University in 1954 under the supervision of Cyril Hinshelwood, working on the chemical kinetics of bacterial growth. Brenner joined the Medical Research Council (MRC) Unit at the Laboratory of Molecular Biology in Cambridge, England, in 1957, where he became one of the most influential molecular biologists of the twentieth century.

H. Robert (Bob) Horvitz was born on May 8, 1947, in Indianapolis, Indiana. He studied mathematics at MIT before switching to biology, earning his Ph.D.\ from Harvard University in 1974 under the supervision of James Watson, working on the molecular biology of bacteriophage. Horvitz then joined the MRC Laboratory of Molecular Biology in Cambridge, where he conducted postdoctoral work with Brenner on the genetics of programmed cell death in the nematode Caenorhabditis elegans. He joined the faculty at MIT in 1978, where he became professor of biology and a Howard Hughes Medical Institute investigator.

John Edward Sulston was born on March 27, 1942, in Buckinghamshire, England. He studied chemistry at Cambridge University and then earned his Ph.D.\ from Cambridge in 1966 under the supervision of Max Perutz, working on the structure of nucleic acids using X-ray crystallography. Sulston joined Brenner's group at the MRC Laboratory of Molecular Biology in Cambridge, where he turned his attention to the biology of C.\ elegans.

The intellectual context of their work was Brenner's decision, in the mid-1960s, to establish a new model organism for studying the genetics of development. Brenner recognized that the powerful genetic tools available for studying bacteria and bacteriophage could not easily be applied to the problem of how genes control the development of a multicellular organism. He needed a simple, multicellular animal with a small number of cells, a short generation time, and a transparent body that could be easily observed. The nematode C.\ elegans, a soil-dwelling roundworm approximately 1 millimeter in length, met all of these criteria.

\section{The Discovery/Contribution}

\subsection{Sydney Brenner and the Establishment of C.\ elegans as a Model Organism}

Brenner's first contribution was the establishment of C.\ elegans as a model organism for the genetic analysis of development. Beginning in the late 1960s, Brenner and his colleagues developed the genetic and molecular tools necessary for working with C.\ elegans: they established mutagenesis protocols, identified thousands of mutant strains, and mapped genes to specific chromosomal locations. They also initiated the painstaking process of constructing a complete cell lineage---a diagram showing the fate of every cell in the organism from the fertilized egg to the adult.

Brenner's choice of C.\ elegans was inspired. The worm has exactly 959 somatic cells in the adult hermaphrodite (and 1031 in the male), and the complete cell lineage was eventually determined by Sulston and others. This invariant cell lineage meant that the fate of each cell was precisely specified, and that genetic mutations affecting cell fate could be identified by their effects on the lineage.

\subsection{John Sulston and the Cell Lineage}

Sulston's contribution was the determination of the complete cell lineage of C.\ elegans. Using a differential interference contrast microscope, Sulston followed individual cells through successive divisions, tracking their fates from the fertilized egg to the L4 larval stage. This work, published in a landmark paper in 1976, revealed that the pattern of cell divisions in C.\ elegans is essentially invariant: every wild-type individual undergoes the same sequence of cell divisions, producing the same cells in the same positions.

Sulston's cell lineage analysis revealed something remarkable: a precise pattern of programmed cell death. Of the 1090 cells produced during the development of the C.\ elegans hermaphrodite, 131 cells die at specific times and in specific locations. These cell deaths are not accidental or pathological---they are part of the normal developmental program and are essential for the formation of the correct body plan.

Sulston also made fundamental contributions to the C.\ elegans genome project, leading the effort to determine the complete DNA sequence of the C.\ elegans genome---the first genome of a multicellular organism to be sequenced, published in 1998.

\subsection{Robert Horvitz and the Genes of Apoptosis}

Horvitz's contribution was the identification and characterization of the genes that control programmed cell death in C.\ elegans. Through a systematic genetic screen, Horvitz and his colleagues identified mutations that either prevented cell deaths that normally occur ( ``cell survival'' mutations) or caused cells to die that normally survive (``cell death'' mutations).

The key genes identified by Horvitz included ced-3, ced-4, and ced-9. The ced-3 gene encodes a protease (now known to be a caspase) that is essential for the execution of cell death. The ced-4 gene encodes an adaptor protein that activates ced-3. Together, ced-3 and ced-4 constitute the core ``death machinery'' of the cell. The ced-9 gene encodes a protein that inhibits ced-3 and ced-4, preventing inappropriate cell death. The balance between ced-9 (the ``survival'' signal) and ced-3/ced-4 (the ``death'' signal) determines whether a cell lives or dies.

Horvitz also identified the genes responsible for the engulfment and disposal of dead cells, including ced-1, ced-2, ced-5, ced-6, ced-7, and ced-10. These genes encode proteins that mediate the recognition and phagocytosis of apoptotic cells by neighboring cells---a process that is essential for the clean removal of dead cells without triggering an inflammatory response.

\section{Impact on Modern Medicine}

The discoveries of Brenner, Horvitz, and Sulston have had a significant impact on our understanding of development, disease, and therapy.

\subsection{Cancer}

Cancer is fundamentally a disease of dysregulated cell death. Many oncogenes promote cell survival by inhibiting apoptosis, while many tumor suppressor genes promote apoptosis. The human homologs of the C.\ elegans death genes---including caspases, Apaf-1, and Bcl-2---play critical roles in cancer. The Bcl-2 protein, which is overexpressed in many cancers, is the human homolog of CED-9 and inhibits apoptosis. Drugs that inhibit Bcl-2---including venetoclax, approved for the treatment of chronic lymphocytic leukemia---restore the apoptotic machinery and kill cancer cells.

\subsection{Autoimmune Disease and Immune Regulation}

Defects in apoptosis are implicated in autoimmune diseases, including systemic lupus erythematosus and rheumatoid arthritis. Inappropriate survival of self-reactive lymphocytes, which should be eliminated by apoptosis during immune cell development, can lead to the production of autoantibodies that attack the body's own tissues. Understanding the genetic mechanisms of apoptosis has guided the development of therapies for autoimmune conditions.

\subsection{Neurodegenerative Disease}

Excessive apoptosis contributes to the neuronal loss that characterizes neurodegenerative diseases, including Alzheimer's disease, Parkinson's disease, and amyotrophic lateral sclerosis. Strategies to inhibit apoptosis in neurons are being explored as potential treatments for these conditions.

\subsection{Developmental Disorders}

Disruptions in the apoptotic program during embryonic development can cause birth defects. The sculpting of tissues and organs---including the fingers, the nervous system, and the cardiovascular system---depends on precisely regulated cell death. Understanding the genetic mechanisms of developmental apoptosis has illuminated the causes of many congenital anomalies.

\section{Legacy}

Sydney Brenner continued his work at the MRC Laboratory of Molecular Biology until 1996, when he moved to the Molecular Sciences Institute in Berkeley, California. He received numerous honors, including the Royal Medal, the Copley Medal, and the Order of Merit. Brenner passed on February 5, 2019, at the age of 91. His vision of using a simple organism to answer fundamental biological questions---and his insistence on choosing the right model organism for the right question---has inspired generations of biologists.

Robert Horvitz continued his work at MIT, where he became David H. Koch Professor of Biology and a Howard Hughes Medical Institute investigator. He received the Wolf Prize in Medicine in 2002 and numerous other honors. Horvitz continues to study the genetic mechanisms of cell death and has been a mentor to many scientists who have gone on to make major contributions to the field.

John Sulston became director of the Wellcome Trust Sanger Institute, where he led the effort to sequence the human genome. He received the Royal Medal and the Order of the British Empire. Sulston was a passionate advocate for open access to scientific data and for the equitable sharing of the benefits of genomic research. He passed on March 6, 2018, at the age of 75.

Together, Brenner, Horvitz, and Sulston revealed that cell death is not a passive event but an active, genetically controlled program that is as important as cell birth in shaping the organism. Their work has fundamentally changed our understanding of development, disease, and the balance between life and death at the cellular level.


\section{Modern Developments}

Brenner, Horvitz, and Sulston's apoptosis work has led to modern cancer therapy. Understanding of programmed cell death has led to BH3 mimetics (venetoclax) for leukemia. Understanding of caspase signaling has informed development of caspase inhibitors for inflammatory diseases. Understanding of apoptosis in development has led to research on neurodegenerative diseases. Understanding of apoptosis in immunity has informed development of immunotherapies.


\section{Bibliography}

\begin{thebibliography}{9}
\bibitem{nobel-2002}
Nobel Prize Outreach, ``The Nobel Prize in Physiology or Medicine 2002,''\emph{NobelPrize.org}.\\
\url{https://www.nobelprize.org/prizes/medicine/2002/summary/}

\bibitem{lecture-2002}
Sydney Brenner, ``Nobel Lecture,''\emph{NobelPrize.org}. Winner-authored primary source; downloadable from the official Nobel Prize archive.\\
\url{https://www.nobelprize.org/prizes/medicine/2002/brenner/lecture/}
\end{thebibliography}
